% ================================================================
% ==== FILE: content/cycles/cycles_k9.tex
% ================================================================

\subsubsection{Zyklen auf \texorpdfstring{$K_9$}{K_9}}
\label{sec:cycles-k9}

Das Niveau$K_9$stellt das metatheoretische Kontinuum dar: den Bereich von
Theorien, Paradigmen, Logiken, Ontologien und methodische Rahmenbedingungen.
Zyklen weiter$K_9$Beschreiben Sie die wiederkehrenden strukturellen Prozesse, durch die
Wissenschaftliche Erkenntnisse entwickeln sich weiter, stabilisieren sich, organisieren sich neu und brechen manchmal zusammen.

Der Zustandsraum ist
\[
  \Omega(K_9)
  =
  \{\text{theories, logics, paradigms, ontologies, models, interpretations}\}.
\]

Die Äxte$A^9(t)$enthalten:
\[
  \begin{aligned}
  A^9 = \{&
     \text{logical distinctions},\;
     \text{ontological distinctions},\\
     &\text{methodological distinctions},\;
     \text{interpretative dimensions}
  \}.
  \end{aligned}
\]

Potenziale$P^9(t)$:
\[
P^9(t)
  = \{\text{explanatory strength},\; \text{predictive power},\;
        \text{simplicity},\; \text{systematicity},\; \text{coherence}\}.
\]

Fließt$J^9(t)$:
\[
  J^9(t) = \{
    \text{argumentation},\;
    \text{proofs},\;
    \text{interpretations},\;
    \text{criticism},\;
    \text{theory transitions}
  \}.
\]

Kritische Schwellenwerte:
\[
  \Theta^9_\mathrm{logic},\quad
  \Theta^9_\mathrm{emp},\quad
  \Theta^9_\mathrm{heur},\quad
  \Theta^9_\mathrm{axiom}.
\]

Ein metatheoretisches Kontinuum existiert, wenn
\[
  k(K_9)>0,\qquad
  T_9 < \min(\Theta^9).
\]


\subsubsection{Überblick}

Zyklen weiter$K_9$kodieren die Dynamik der Theoriebildung, Stabilisierung,
Bewertung, Korrektur und revolutionärer Ersatz. Sie formalisieren
die systemischen Prozesse, die dem wissenschaftlichen Wissen und seiner Entwicklung zugrunde liegen.
These cycles operate on the structural space $\Omega(K_9)$ and are driven
durch logische, empirische und methodische Zwänge, die durch ausgedrückt werden$T_9$.


\subsubsection{1. Der normale Wissenschaftszyklus$C_{\mathrm{norm}}$}

Dieser Zyklus entspricht stabilen Perioden theoriegeleiteter Forschung.

\[
  \begin{aligned}
  \mathrm{Puzzle\;selection}
  &\to \mathrm{Method\;application}
  \to \mathrm{Result\;integration}\\
  &\to \mathrm{Coherence\;increase}
  \to \mathrm{Puzzle\;selection}.
  \end{aligned}
\]

Formal:
\[
 C_{\mathrm{norm}}
 =
 \left\{
   J^9_{\mathrm{arg}},\;
   P^9_{\mathrm{coherence}},\;
   P^9_{\mathrm{systematicity}}
 \right\},
\]
where $J^9_{\mathrm{arg}}$ is the argumentative flow supporting integration.

Stabilität erfordert:
\[
  \Theta^9_\mathrm{logic} < P^9_{\mathrm{coherence}}.
\]


\subsubsection{2. Der Beweis-Widerlegungs-Zyklus$C_{\mathrm{proof/ref}}$}

Der grundlegende logische Zyklus in der Wissenschaft.

\[
  \begin{aligned}
  \mathrm{Hypothesis}
  &\to \mathrm{Derivation}
  \to \mathrm{Test}
  \to \mathrm{Refutation\;or\;support}\\
  &\to \mathrm{Refinement}
  \to \mathrm{Hypothesis}.
  \end{aligned}
\]

Bilden:
\[
 C_{\mathrm{proof/ref}}
  =
  \left\{
    J^9_{\mathrm{proof}},\;
    J^9_{\mathrm{crit}},\;
    P^9_{\mathrm{predictive}},\;
    \Theta^9_{\mathrm{emp}}
  \right\}.
\]

Schwelle:
\[
  P^9_{\mathrm{predictive}} > \Theta^9_\mathrm{emp}.
\]


\subsubsection{3. Der Interpretationszyklus$C_{\mathrm{interp}}$}

Interpretationszyklen reorganisieren Theorien, ohne den empirischen Inhalt zu verändern.

\[
  \mathrm{Reformulation}
  \to
  \mathrm{Conceptual\;mapping}
  \to
  \mathrm{Reduction\;or\;extension}
  \to
  \mathrm{Unification}
  \to
  \mathrm{Reformulation}.
\]

Bilden:
\[
 C_{\mathrm{interp}}
  =
  \left\{
    A^9_{\mathrm{interpretation}},\;
    P^9_{\mathrm{simplicity}},\;
    J^9_{\mathrm{reinterpretation}}
  \right\}.
\]

Stabilitätsbedingung:
\[
  P^9_{\mathrm{simplicity}} >
  \Theta^9_{\mathrm{heur}}.
\]


\subsubsection{4. Der Paradigmenzyklus$C_{\mathrm{paradigm}}$}

Dieser Zyklus beschreibt die langfristige Entwicklung ganzer wissenschaftlicher Paradigmen.

\[
  \begin{aligned}
  \mathrm{Model\;development}
  &\to \mathrm{Problem\;solving}
  \to \mathrm{Anomaly\;accumulation}\\
  &\to \mathrm{Crisis}
  \to \mathrm{Reconstruction}
  \to \mathrm{Model\;development}.
  \end{aligned}
\]

Bilden:
\[
 C_{\mathrm{paradigm}}
  =
  \left\{
    P^9_{\mathrm{explanatory}},\;
    J^9_{\mathrm{crit}},\;
    A^9_{\mathrm{ontology}},\;
    \Theta^9_{\mathrm{axiom}}
  \right\}.
\]

Schwelle (Kuhn-Typ):
\[
  T_9 > \Theta^9_\mathrm{logic}
  \quad\Rightarrow\quad
  \text{Paradigm instability}.
\]


\subsubsection{5. Der Zyklus der wissenschaftlichen Revolution$C_{\mathrm{rev}}$}

Der schnelle, spannungsgeladene Übergang, bei dem ein neues Paradigma das alte ersetzt.

\[
  \mathrm{Crisis}
  \to
  \mathrm{Competing\;frameworks}
  \to
  \mathrm{Selection}
  \to
  \mathrm{Reinterpretation\;of\;past}
  \to
  \mathrm{Stabilisation}
  \to
  \mathrm{Crisis}.
\]

Formal:
\[
 C_{\mathrm{rev}}
   =
   \left\{
     J^9_{\mathrm{transition}},\;
     A^9_{\mathrm{methodology}},\;
     P^9_{\mathrm{explanatory}},\;
     P^9_{\mathrm{predictive}}
   \right\}.
\]

Die Übergangsbedingung (metatheoretischer Phasenübergang):
\[
  T_9 = \Theta^9_{\mathrm{dim}}.
\]

Cycle reinstates a new structure of $\Omega(K_9)$.


\subsubsection{6. Der logische Konsistenzzyklus$C_{\mathrm{logic}}$}

Basierend auf interner Logik und Axiomatik.

\[
  \mathrm{Axiom\;selection}
  \to
  \mathrm{Derivation}
  \to
  \mathrm{Consistency\;check}
  \to
  \mathrm{Revision}
  \to
  \mathrm{Axiom\;selection}.
\]

Bilden:
\[
 C_{\mathrm{logic}}
  =
  \left\{
    A^9_{\mathrm{logic}},\;
    \Theta^9_{\mathrm{axiom}},\;
    J^9_{\mathrm{proof}}
  \right\}.
\]

Verbunden mit Grenzwerten vom Typ Gödel:
\[
  P^9_{\mathrm{coherence}} > \Theta^9_{\mathrm{axiom}}
  \;\;\text{for viability}.
\]


\subsubsection{7. Der Evolutionszyklus der Metatheorie$C_{\mathrm{meta}}$}

Dieser Zyklus regelt Übergänge zwischen gesamten Meta-Frameworks und Logiken.

\[
  \mathrm{Framework}
  \to
  \mathrm{Cross\!-\!comparison}
  \to
  \mathrm{Meta\!-\!critique}
  \to
  \mathrm{Higher\!-\!order\;reconstruction}
  \to
  \mathrm{Framework}.
\]

Bilden:
\[
 C_{\mathrm{meta}}
 =
 \left\{
   A^9_{\mathrm{paradigm}},\;
   A^9_{\mathrm{methodology}},\;
   J^9_{\mathrm{meta}},\;
   \Theta^9_{\mathrm{heur}}
 \right\}.
\]

Zustand:
\[
  k(K_9) > 0 \quad\Leftrightarrow\quad
  C_{\mathrm{meta}}\;\text{remains stable}.
\]


\subsubsection{Zyklusmetriken für \texorpdfstring{$K_9$}{K_9}}

Using the general definitions (the corresponding derivation note):

\paragraph{Cycle length.}
\[
  L(C) = \oint d\Omega(K_9).
\]

\paragraph{Efficiency.}
\[
  C_{\mathrm{eff}} =
  \frac{\oint J^9 \cdot dA^9}{L(C)}.
\]

\paragraph{Stability.}
\[
SC)
    = \min_{t\in C}
      d(\Omega(K_9),\partial\Omega(K_9)).
\]

\paragraph{Weight.}
\[
  w(C) = F(C_{\mathrm{eff}}, S(C),
          P^9_{\mathrm{coherence}},
          P^9_{\mathrm{predictive}}).
\]


\subsubsection{Zusammenbruch von \texorpdfstring{$K_9$}{K_9} Zyklen}

Ein Kollaps tritt auf, wenn:

\[
  P^9_{\mathrm{coherence}} < \Theta^9_{\mathrm{logic}},\qquad
  P^9_{\mathrm{predictive}} < \Theta^9_{\mathrm{emp}},\qquad
  P^9_{\mathrm{simplicity}} < \Theta^9_{\mathrm{heur}}.
\]

Dann:

\[
  C_j \to \varnothing,\qquad
  \Omega(K_9)=\varnothing,\qquad
  k(K_9)\to 0.
\]

Dies entspricht dem Tod eines Paradigmas oder Meta-Frameworks.


\subsubsection{Kontinuität von \texorpdfstring{$K_8$}{K_8} zu \texorpdfstring{$K_9$}{K_9}}

The transition $K_8\to K_9$ occurs when:

\[
  T_8 > \Theta^9_{\mathrm{dim}},
\]

und neue reflektierende Achsen erscheinen:
\[
  \begin{aligned}
  A^9_{\mathrm{logic}},\;
  A^9_{\mathrm{ontology}},\\
  A^9_{\mathrm{methodology}},\;
  A^9_{\mathrm{interpretation}}.
  \end{aligned}
\]

Dadurch entsteht ein neues metatheoretisches Kontinuum mit eigenen Zyklen und Stabilität
Bedingungen.
